\documentclass{situationist-preprint}

\title{Phenomenal Concepts After Release}
\msid{MND-04-b}
\msfamily{MND-04}
\msdiscipline{Philosophy of Mind}
\mstype{Consequence}
\mskeywords{Mary's room; phenomenal concepts; ability hypothesis; a posteriori identity; explanatory gap}

\begin{document}
\maketitle

\begin{abstract}
MND-04 argues that Mary, on leaving her room, learns no new fact but acquires a recognitional ability and with it a phenomenal concept of a fact she already knew. We derive three consequences for the concept she acquires. It co-refers with a physical concept she already had; the identity between the two is knowable only a posteriori; and her new position explains, rather than contradicts, her continuing sense that the physical facts leave something out. The consequences connect the result of MND-04 with the account of zombies in MND-05 and of the explanatory gap in MND-01.
\end{abstract}

\section{The concept Mary acquires}

According to Theorem~3.3 of MND-04-a, what Mary acquires when she first sees a ripe tomato is a recognitional ability, the ability to tell by having an experience that it is of a given type (Definition~2.1), and a phenomenal concept, a concept whose possession consists in or requires that ability (Definition~2.2). The concept is of a fact she already knew: by Lemma~3.2, she knows the same facts about colour vision before and after her release. Call the type of experience she has on seeing red $E$, the physical concept under which she knew of $E$ in the room $P$, and the phenomenal concept she acquires $R$.

\begin{proposition}\label{p:coref}
After her release, Mary possesses two concepts, $P$ and $R$, that refer to the same type of experience.
\end{proposition}

\begin{proof}
$R$ is a concept of the type of experience she recognizes, which is $E$. In the room she knew every physical fact about colour vision, and in particular knew of $E$ under $P$. By Lemma~3.2 of MND-04-a, the fact that $R$ lets her think about is one she already knew; so what $R$ refers to is what $P$ refers to.
\end{proof}

\begin{proposition}\label{p:apost}
The identity $R = P$ is not knowable a priori, even by Mary after her release.
\end{proposition}

\begin{proof}
A phenomenal concept refers to its object directly, through a recognitional capacity, without a descriptive mode of presentation that could reveal its connection to a physical concept (MND-05-a, Definition~2.1). So no analysis of $R$ yields the physical description that $P$ expresses, and no analysis of $P$ yields the recognitional capacity on which $R$ depends; this is the conceptual isolation of Assumption~2.2 of MND-05-a. Mary can come to know the identity, but only by an a posteriori route: by finding that the recognitional ability is triggered by the very neural states that $P$ picks out.
\end{proof}

The identity is like that of water with H$_2$O in one respect and unlike it in another. Like it, it is necessary, if true, and known a posteriori. Unlike it, it is not mediated by a contingent description such as \enquote{the clear liquid in lakes}, under which water is picked out and which H$_2$O could be found to satisfy. The phenomenal concept presents its referent directly \citep{Loar1990}. This is why the identity, unlike that of water, continues to seem contingent after it is known.

\section{The residual sense of a gap}

\begin{corollary}\label{c:gap}
After her release, Mary will find the physical facts insufficient to tell her, a priori, which experiences she is having, although she lacks no fact.
\end{corollary}

\begin{proof}
By Proposition~\ref{p:apost}, she cannot derive from the physical facts expressed under $P$ the truth she expresses under $R$ without an a posteriori premise. By Proposition~\ref{p:coref} and Lemma~3.2 of MND-04-a, the truth she expresses under $R$ is a fact she already knew. Her sense that the physical facts leave something out is therefore a sense of a gap between concepts, not between facts.
\end{proof}

Corollary~\ref{c:gap} is the form, in Mary's case, of the explanation of the explanatory gap given by MND-01: the same facts are presented in two ways, structurally in physics and intrinsically in experience, and no a priori derivation runs from the first to the second. It is also the reason the conceivability of zombies, which MND-05 grants, does not establish their possibility: the propositions that make zombies conceivable are formulated with concepts such as $R$.

\begin{objection}
On the ability hypothesis, Mary gains know-how and not knowledge-that \citep{Lewis1988,Nemirow1990}. To credit her with a new concept is to concede that she gains new knowledge-that, and so a new fact.
\end{objection}

\begin{reply}
MND-04 credits her with both, and the consequence derived here shows why the concession costs nothing. New knowledge-that, in the sense of a new way of thinking of a fact, is not knowledge of a new fact. The knowledge argument \citep{Jackson1982} is valid only if the first entails the second, and Theorem~3.3 of MND-04-a shows that it does not. The ability is what grounds the concept; the concept is what lets her entertain a known fact in a new way.
\end{reply}

\section{Conclusion}

After her release Mary has two concepts of the same experience, one physical and one phenomenal. They co-refer, their identity is knowable only a posteriori, and the phenomenal concept's directness explains why the identity continues to seem contingent. Her new position thus reproduces, in a single subject, the explanatory gap that MND-01 explains and the conceivability of zombies that MND-05 grants, without any new fact.

\begin{thebibliography}{9}\small
\bibitem[Jackson(1982)]{Jackson1982} Jackson, F. (1982). Epiphenomenal qualia. \emph{The Philosophical Quarterly}, 32(127), 127--136.
\bibitem[Lewis(1988)]{Lewis1988} Lewis, D. (1988). What experience teaches. \emph{Proceedings of the Russellian Society}, 13, 29--57.
\bibitem[Loar(1990)]{Loar1990} Loar, B. (1990). Phenomenal states. \emph{Philosophical Perspectives}, 4, 81--108.
\bibitem[Nemirow(1990)]{Nemirow1990} Nemirow, L. (1990). Physicalism and the cognitive role of acquaintance. In W.~G. Lycan (ed.), \emph{Mind and Cognition: A Reader}. Blackwell.
\end{thebibliography}

\end{document}
